\chapter{2025年上海卷（秋）}

\section{填空题}

\begin{problem}
已知全集 \(U = \{x \mid 2 \le x \le 5, x \in \R\}\)，集合 \(A = \{x \mid 2 \le x < 4, x \in \R\}\)，则 \(\bar{A} =\)\fillinblank{}.
\end{problem}

\begin{problem}
设 \(x\in\R\)，不等式 \(\frac{x-1}{x-3}<0\) 的解集为\fillinblank{}.
\end{problem}

\begin{problem}
已知等差数列 \(\{a_n\}\) 的首项 \(a_1 = -3\)，公差 \(d = 2\)，则该数列的前 \(6\) 项和为\fillinblank{}.
\end{problem}

\begin{problem}
在二项式 \((2x-1)^5\) 的展开式中，\(x^3\) 的系数为\fillinblank{}.
\end{problem}

\begin{problem}
函数 \(y=\cos x\) 在 \([-\pi/2,\pi/4]\) 上的值域为\fillinblank{}.
\end{problem}

\begin{problem}
已知随机变量 \(X\) 的分布为
\(\begin{pmatrix}5&6&7\\0.2&0.3&0.5\end{pmatrix}\)，则期望 \(E(X)=\)\fillinblank{}.
\end{problem}

\begin{problem}
如图，在正四棱柱 \(ABCD-A_1B_1C_1D_1\) 中，\(AC=4\sqrt2\)，\(BD_1=9\)，则该正四棱柱的体积为\fillinblank{}.

\bitmapfigure[max width=6.0cm]{img/2025/shanghai/q7_fig1.png}
\end{problem}

\begin{problem}
设 \(a\)，\(b>0\)，若 \(a+\frac1b=1\)，则 \(b+\frac1a\) 的最小值为\fillinblank{}.
\end{problem}

\begin{problem}
\(4\) 个家长和 \(2\) 个儿童去郊游爬山，\(6\) 个人需要排成一条队列，要求队列的头和尾均是家长，则不同的排列个数有\fillinblank{}.
\end{problem}

\begin{problem}
\(i\) 为虚数单位，若复数 \(z\) 满足 \(z^2=(\bar z)^2\)，且 \(|z|\le1\)，则 \(|z-2-3\i|\) 的最小值是\fillinblank{}.
\end{problem}

\begin{problem}
小申同学观察发现，生活中有时候影子可以完全投射在斜面上. 某斜面上有两根长为 \(1\) 米且垂直于水平面放置的杆子，与斜面的接触点分别为 \(A\)，\(B\)，它们在阳光的照射下呈现出影子，阳光可视为平行光，其中 \(A\) 处杆子的影子在水平面上，长度为 \(0.4\) 米，\(B\) 处杆子的影子完全在斜面上，长度为 \(0.45\) 米. 则斜面的坡角 \(\theta=\)\fillinblank{}（结果精确到 \(0.01^\circ\)）.

\bitmapfigure[max width=6.0cm]{img/2025/shanghai/q11_fig1.png}
\end{problem}

\begin{problem}
已知函数 \(f(x)\) 满足：当 \(x>0\) 时 \(f(x)=1\)，当 \(x=0\) 时 \(f(x)=0\)，当 \(x<0\) 时 \(f(x)=-1\). \(\bs{a}\)，\(\bs{b}\)，\(\bs{c}\) 分别为平面内三个不同的单位向量. 若 \(f(\bs{a}\cdot\bs{b})+f(\bs{b}\cdot\bs{c})+f(\bs{c}\cdot\bs{a})=0\)，则 \(|\bs{a}+\bs{b}+\bs{c}|\) 的可取值范围是\fillinblank{}.
\end{problem}

\section{单选题}

\begin{problem}
已知事件 \(A\)，\(B\) 相互独立，若 \(P(A)=P(B)=\frac{1}{2}\)，则 \(P(A\cap B)\) 的值为（\quad）

\choices
    {\(0\)}
    {\(\frac{1}{4}\)}
    {\(\frac{1}{2}\)}
    {\(1\)}
\end{problem}

\begin{problem}
已知 \(a>0\) 且 \(a\ne1\)，\(s\in\R\)，下列各项中，能推出 \(a^s>a\) 的是（\quad）

\choices
    {\(a>1\) 且 \(s>0\)}
    {\(a>1\) 且 \(s<0\)}
    {\(0<a<1\) 且 \(s>0\)}
    {\(0<a<1\) 且 \(s<0\)}
\end{problem}

\begin{problem}
已知 \(A(0,1)\)，\(B(1,2)\)，\(C\) 在 \(\Gamma:x^2-y^2=1\ (x\ge1,y\ge0)\) 上，则 \(\triangle ABC\) 的面积（\quad）

\choices
    {有最大值，但没有最小值}
    {没有最大值，但有最小值}
    {既有最大值，也有最小值}
    {既没有最大值，也没有最小值}
\end{problem}

\begin{problem}
设 \(\lambda\in\R\)，数列 \(\{a_n\}\)，\(\{b_n\}\)，\(\{c_n\}\) 的通项分别为 \(a_n=10n-9\)，\(b_n=2^n\)，\(c_n=\lambda a_n+(1-\lambda)b_n\). 若对任意 \(\lambda\in[0,1]\)，\(a_n\)，\(b_n\)，\(c_n\) 均能作为三角形的三边长，则满足条件的正整数 \(n\) 有（\quad）

\choices
    {1个}
    {3个}
    {4个}
    {无穷多个}
\end{problem}

\section{解答题}

\begin{problem}
2024 年东京奥运会，中国获得了男子 \(4\times100\) 米混合泳接力金牌，以下是历届奥运会男子 \(4\times100\) 米混合泳接力项目冠军成绩（单位：秒），数据按照升序排列.

\begin{center}
\begin{tabular}{|c|c|c|c|c|}
\hline
206.78 & 207.46 & 207.95 & 209.34 & 209.35 \\
\hline
210.68 & 213.73 & 214.84 & 216.93 & 216.93 \\
\hline
\end{tabular}
\end{center}

\begin{enumerate}
\item 求这组数据的极差与中位数；
\item 从这 10 个数据中任选 3 个，求恰有 2 个数据大于 \(211\) 的概率；
\item 若比赛成绩 \(y\) 关于年份 \(x\) 的回归方程为 \(y=-0.311x+\hat b\)，年份 \(x\) 的平均数为 \(2006\)，预测 \(2028\) 年男子 \(4\times100\) 米混合泳接力项目冠军队的成绩（精确到 \(0.01\) 秒）.
\end{enumerate}
\end{problem}

\begin{problem}
如图，\(P\) 是圆锥顶点，\(O\) 为底面圆心，\(AB\) 为底面直径，且 \(AB=2\).
\begin{enumerate}
\item 若直线 \(PA\) 与圆锥底面所成角为 \(\pi/3\)，求圆锥侧面积；
\item 已知 \(Q\) 是母线 \(PA\) 中点，点 \(C\)，\(D\) 在底面圆周上，且弧 \(AC\) 长为 \(\pi/3\)，\(CD\parallel AB\). 设点 \(T\) 在线段 \(OC\) 上，证明 \(QT\parallel PBD\) 所在平面.
\end{enumerate}

\bitmapfigure[max width=6.0cm]{img/2025/shanghai/q18_fig1.png}
\end{problem}

\begin{problem}
已知 \(f(x)=x^2-(m+2)x+m\ln x\)，\(m\in\R\).
\begin{enumerate}
\item 若 \(f(1)=0\)，求不等式 \(f(x)\le x^2-1\) 的解集；
\item 若 \(f(x)\) 在 \((0,+\infty)\) 上存在极大值，求 \(m\) 的取值范围.
\end{enumerate}
\end{problem}

\begin{problem}
已知椭圆 \(\Gamma:\frac{x^2}{a^2}+\frac{y^2}{5}=1\)（\(a>\sqrt5\)），\(M(0,m)\)（\(m>0\)），\(A\) 为 \(\Gamma\) 的右顶点.
\begin{enumerate}
\item 若 \(\Gamma\) 的焦点为 \((2,0)\)，求离心率 \(e\)；
\item 若 \(a=4\)，且 \(\Gamma\) 上存在一点 \(P\) 满足 \(\overrightarrow{PA}=2\overrightarrow{MP}\)，求 \(m\)；
\item 若线段 \(AM\) 的中垂线 \(l\) 的斜率为2， \(l\) 与 \(\Gamma\) 交于 \(C\)，\(D\)，且 \(\angle CMD\) 为钝角，求 \(a\) 的取值范围.
\end{enumerate}
\end{problem}

\begin{problem}
已知函数 \(y=f(x)\) 定义域为 \(\R\). 对于正实数 \(a\)，定义集合 \(M_a=\{x\mid f(x+a)=f(x)\}\).
\begin{enumerate}
\item 若 \(f(x)=\sin x\)，判断 \(\dfrac{\pi}{3}\) 是否是 \(M_{\pi}\) 的元素，请说明理由；
\item 若 \(f(x)=x+2\)（\(x<0\)），\(f(0)=0\)，且当 \(x>0\) 时 \(f(x)=\sqrt{x}\)，\(M_a\neq\varnothing\)，求 \(a\) 的取值范围；
\item 若 \(y=f(x)\) 是偶函数，当 \(x\in(0,1]\) 时 \(f(x)=1-x\)，且对任意 \(a\in(0,2)\)，有 \(M_a\subset M_2\). 写出 \(f(x)\) 在 \(x\in(1,2)\) 的解析式，并证明：对任意实数 \(c\)，函数 \(f(x)-c\) 在 \([-3,3]\) 上至多有9个零点.
\end{enumerate}
\end{problem}

